\documentclass[11pt]{article}
\usepackage[T1]{fontenc}
\usepackage[utf8]{inputenc}
\usepackage{lmodern}
\usepackage[margin=2.5cm]{geometry}
\usepackage{amsmath,amssymb,amsthm}
\usepackage{enumitem}
\usepackage{microtype}
\usepackage[colorlinks=true,linkcolor=blue,urlcolor=blue]{hyperref}
\usepackage{xurl}
\usepackage{listings,xcolor}
\lstset{basicstyle=\ttfamily\small,breaklines=true,columns=fullflexible,frame=single,keepspaces=true,upquote=true,literate={≤}{{$\le$}}1 {≥}{{$\ge$}}1 {−}{{-}}1 {×}{{$\times$}}1}
\providecommand{\tightlist}{\setlength{\itemsep}{0pt}\setlength{\parskip}{0pt}}
\setlength{\parindent}{0pt}
\setlength{\parskip}{0.6em}
\allowdisplaybreaks

\title{Almost-linear-time lossy min-cost flow with prescribed demands}
\date{}

\begin{document}
\maketitle

Let \(A_0\) be the imbalance matrix of the directed graph. A column for \(e=(v,w)\) is \(-\mathbf e_v+\gamma_e\mathbf e_w\), where \(\mathbf e_v\) is the vertex unit vector and \(0<\gamma_e\le1\). On a loop this is the single entry \(\gamma_e-1\). Write \(m_{\rm in}\) for the number of supplied arcs and \(m\) for the number with positive capacity. For the end-to-end bound use \(N=m_{\rm in}+n+2\). If size is measured after filtering, the bound uses \(N=m+n+2\) on the filtered input, with the scan of any unfiltered list and expansion of discarded output positions charged as input and output work. All extrema in \(W=\max(1,\max\gamma_e^{-1},\|c\|_\infty,\max u_e,\max u_e/\min u_e)\) are over positive-capacity arcs. Assume \(A_0x=d\) is attainable in the capacity box. For each \(\delta>0\) the algorithm below returns
\[
 0\le f\le u,\qquad \|A_0 f-d\|_\infty\le\delta,\qquad
 c^\top f\le \min_{A_0 x=d,\ 0\le x\le u}c^\top x .
\]
With high probability the work is \(N^{1+o(1)}\operatorname{polylog}(2+NW/\delta)\) in the precision-access model specified below. The algorithm solves the bounded intermediate problem in polylogarithmically many rounds using ordinary integral min-cost flow as its specialized algorithmic primitive. A dyadic is an integer times an integer power of two. A mesh \(h\), always a power of two, denotes the grid \(h\mathbb Z\). Rounding down, up, or to the nearest grid value means taking the corresponding multiple of \(h\), with ties for nearest rounding resolved downward. A box \([l,d]\) has coordinatewise endpoints, and scalar bounds on vectors apply to every coordinate.

\hypertarget{an-ordinary-flow-minimizer-for-homogeneous-directed-hinges}{%
\section*{An ordinary-flow minimizer for homogeneous directed hinges}\label{an-ordinary-flow-minimizer-for-homogeneous-directed-hinges}}

Use a single physical root pinned at one to represent unary hinges. Its relative coordinate is pinned at zero. On a box of positive vertex values consider
\[
 F(x)=\ell^\top x+\sum_e k_e[a_e x_i-b_e x_j]_+ +\mathrm{constant},
 \quad [t]_+=\max(t,0),\quad a_e,b_e,k_e\ge0 .
\]
For a convex piecewise-linear unary, put its first slope in \(\ell\) and represent a slope increase \(k\) at a positive breakpoint \(t\) by \(k[x_i-tx_0]_+\) with \(x_0=1\). Use \(\wedge,\vee\) for coordinatewise minimum and maximum. Each hinge term \(F_e(x)=k_e[a_ex_i-b_ex_j]_+\) is convex and satisfies \(F_e(x\wedge y)+F_e(x\vee y)\le F_e(x)+F_e(y)\). To see the latter for distinct endpoints, take \(s\le s'\) in the first coordinate and \(t\le t'\) in the second coordinate. The increase from \(s\) to \(s'\) at second coordinate \(t\) is \(k_e([r+\Delta]_+-[r]_+)\), with \(r=a_es-b_et\) and \(\Delta=a_e(s'-s)\ge0\). The expression is nondecreasing in \(r\), so the increase at \(t\) is at least the increase at \(t'\). Rearranging gives the inequality for the crossed pair \((s,t'),(s',t)\). For a coordinatewise ordered pair the two sides agree. A pinned endpoint or a common endpoint yields a unary function, which has equality in the lattice inequality. Linearity and addition now give the inequality for \(F\).

Suppose the objective range is at most \(V\) on the original box \([l,d^{(0)}]\), and write \(F^*\) for its minimum, attained by continuity on the compact box. A box with no free coordinates returns its single point. Starting at \(d=d^{(0)}\), minimize to additive \(\eta\) in \([\max(l,d/2),d]\) and use the output \(v\) as the next upper vector. For analysis maintain a point \(u\in[l,d]\) of score at most \(F^*+j\eta\) after \(j\) rounds. Initially \(u\) is an exact optimum. Since \(u\vee v\) is in the current box, \(F(v)\le F(u\vee v)+\eta\). The lattice inequality gives \(F(u\wedge v)\le F(u)+F(v)-F(u\vee v)\le F(u)+\eta\), so the next analysis point can be \(u\wedge v\). The midpoint of \(d,u\) belongs to the current box, hence convexity gives
\[
 F(v)-F^*\le (F(d)-F^*)/2+(j/2+1)\eta .
\]
For \(G_j=F(d^{(j)})-F^*\) this gives \(G_{j+1}\le G_j/2+(j/2+1)\eta\). The bound \(G_j\le 2^{-j}V+(j+1)\eta\) holds at \(j=0\), and substitution bounds \(G_{j+1}\) by \(2^{-j-1}V+(j+3/2)\eta\). With \(T_0=\lceil\log_2(4+4V/\xi)\rceil\) and \(\eta\) the largest power of two at most \(\xi/[8(T_0+2)]\), the two terms after \(T_0\) rounds are at most \(\xi/4\) and \(\xi/8\), respectively.

Start each factor-two minimization at the upper vector of its fixed box. About the current inner point \(d\), write \(x_i=d_i(1+z_i)\), with \(d_0=1,\ z_0=0\) at the root. The ratios of upper to lower bounds of the fixed box are at most two, so its relative intervals lie in \([-1/2,1]\). For one hinge with positive endpoint coefficients put \(P=a_e d_i,\ Q=b_e d_j,\ r=P-Q\). Its new argument is \(A(z)=r(1+z_j)+P(z_i-z_j)\). If \(r>0\), write \([A(z)]_+=A(z)+[-r(1+z_j)-P(z_i-z_j)]_+\). If \(r\le0\), the affine part is zero and the remainder is \([-|r|(1+z_j)+P(z_i-z_j)]_+\). In either remainder replace \(1+z_j\) by \(1/2\). Since \(1+z_j\ge1/2\), this gives an upper model \(\mathcal M\), with equality at \(z=0\). At a quarter of the relative displacement \(z^*\) to a box optimum, the affine change is a quarter of its change at \(z^*\). The modeled remainder is at most a quarter of the true remainder at \(z^*\) by \(1/2\ge(1+z_j^*)/4\) and positive homogeneity of the positive part. Both remainders vanish at zero. The quarter displacement is inside the fixed box, and summing over terms gives \(\mathcal M(z^*/4)\le3F(d)/4+F(x^*)/4\). Each model consists of translated ordinary difference hinges and an affine term. If an endpoint coefficient is zero, the hinge is either zero or linear on positive physical coordinates and is modeled exactly. The procedure below returns a physical point of score at most \(\mathcal M(z^*/4)+\eta/32\). Use just this inner loop when the original box already has coordinatewise ratio at most two.

Each such ordinary model has terms \(k_e[z_i-z_j-c_e]_+\) and a linear term \(b^\top z\). For dyadic data impose the relative interval \([\ell_i,r_i]\) at each nonroot coordinate by adding \(K[z_i-z_0-r_i]_++K[z_0-z_i+\ell_i]_+\), where \(K>1+\|b\|_1+2\sum_e k_e\) before adding these penalties. The absolute unpenalized model change from \(z\) to its clipped value \(\bar z\) is bounded by \((\|b\|_1+2\sum_e k_e)\|z-\bar z\|_\infty\). Each difference in a hinge changes by at most \(2\|z-\bar z\|_\infty\). The penalties saved are exactly \(K\|z-\bar z\|_1\), so clipping any point outside strictly improves its penalized score. Minimizing the continuous model over the compact box now gives an unconstrained minimizer, and every unconstrained minimizer lies in the box.

Let \(B\) have row \(\mathbf e_i^\top-\mathbf e_j^\top\) for the hinge at \(i,j\), and give \(b\) the root entry \(b_0=-\sum_{i\ne0} b_i\). Use \(c,k\) for the vectors of shifts and slopes including the interval penalties. With \(z_0=0\), introduce an epigraph variable \(q_e\ge0\) with \(q_e\ge (Bz)_e-c_e\) for each hinge. In the Lagrangian the multiplier \(s_e\ge0\) of \((Bz)_e-c_e-q_e\le0\) leaves the coefficient \(k_e-s_e\) on \(q_e\). Minimization over \(q\ge0\) gives \(s\le k\), and minimization over free \(z_i\) for \(i\ne0\) gives \(B^\top s=-b\) there. The root equation follows since \(B\mathbf1=0\). Strong LP duality gives the value \(\max_{B^\top s=-b,\ 0\le s\le k}(-c^\top s)\). Orient arc \(e\) from \(i\) to \(j\), so \(B^\top s\) is net outflow and the ordinary min-cost flow uses net inflow demand \(b\) and cost \(c\). Route \(|b_i|\) on the interval arc \(i\to0\) if \(b_i<0\), on \(0\to i\) if \(b_i>0\), and use zero if \(b_i=0\). Each capacity \(K\) suffices and the root's demand follows by balance. Drop zero-slope and constant self-hinges. Scale all capacities and demands by the reciprocal of their common dyadic mesh and scale costs by the reciprocal of their own mesh. The result is an integral min-cost flow instance with an explicit integral feasible flow given by this routing. To use a simple directed graph, subdivide each arc through its own private vertex, giving both pieces the original capacity, the first its original cost and the second zero cost. Give the private vertex zero demand. Its equality makes the two flows identical, so projection preserves all requirements and costs.

To obtain an optimal potential from the projected optimal flow, complementary slackness requires \(z_i-z_j\le c_e\) when \(s_e<k_e\) and \(z_j-z_i\le -c_e\) when \(s_e>0\). These inequalities ensure \(s_e\) maximizes \(v((Bz)_e-c_e)\) over \(0\le v\le k_e\), so the primal hinge model then attains the flow bound. Test the conditions on \(s_e\) using the scaled flow and capacity integers. Let \(\hat c_e\) be the integer obtained by scaling \(c_e\). Insert an arc \(j\to i\) of length \(\hat c_e\) for the first inequality and \(i\to j\) of length \(-\hat c_e\) for the second. There is a scaled potential satisfying all inequalities by strong duality and slackness. Summing its inequalities around a directed cycle shows that the cycle has nonnegative length.

For any such directed integer-length graph without negative cycles, compute distance labels as follows. Add a new source with zero-length arcs to all \(v_o\) old vertices, one unit of net inflow demand at each old vertex, and net supply \(v_o\) at the source. Give every arc capacity \(v_o+1\). An explicit feasible flow sends one unit on each source arc. Every old vertex is reachable, and removing cycles from walks without increasing their length shows that a simple path attains each shortest distance \(\pi_i\). The distances satisfy \(\pi_j-\pi_i\le C_{ij}\) on each arc \(i\to j\) of length \(C_{ij}\). Sending one unit along a simple shortest path to each old vertex respects the capacity \(v_o+1\) and has cost \(\sum_{\rm old}\pi_i\). Summing the distance differences against any feasible routing gives this sum as a lower bound on its cost. Hence every positive arc of an optimal routing has zero reduced length \(C_{ij}+\pi_i-\pi_j\). If some vertex were unreachable from the source in the positive support, summing net inflow over the unreachable set would give positive total demand with no incoming positive flow. A search tree of positive arcs therefore reaches every vertex. Summing lengths along that tree gives \(\pi\). Apply this procedure to the inequality graph, subtract the old root distance, and undo cost scaling to get the desired potential. Use subdivision after assigning the demands and project the returned flow back before the traversal. Computing these labels uses one ordinary solve and one traversal.

For precision, let \(C_0\ge2\) bound the number of model vertices and hinges before adding penalties, all absolute shifts, the sum \(\|b\|_1+2\sum_e k_e\), the largest physical scale, and a maximum-norm Lipschitz constant of the physical objective. Round each relative interval outward by at most \(h\). Round each shift, slope and linear coefficient to the nearest multiple of \(h\), preserving nonnegative slopes, where \(h\le\min(1/2,\eta/(10^5 C_0^4))\). The rounded relative box lies in \(|z_i|\le2\). On this box a hinge with rounded slope \(k'_e\) and shift \(c'_e\) has value error at most \(h(4+C_0)+k'_e h\). Indeed, \(|z_i-z_j|\le4\), and changing the shift changes the positive part by at most \(h\). The linear term changes by at most \(2h\) per coefficient, and \(\sum_e k'_e\le \sum_e k_e+C_0h\). Thus the rounded model \(\mathcal M_h\) differs from \(\mathcal M\) by at most \(10C_0^2 h\), with the same additive constant understood in both.

Add the interval penalties to \(\mathcal M_h\) using \(K\) chosen for the rounded slopes. Its optimal potential \(z^{\rm rd}\) lies in the outward-rounded box. Clip it to the true relative box to obtain \(z^{\rm cl}\), with \(\|z^{\rm cl}-z^{\rm rd}\|_\infty\le h\). By optimality, \(\mathcal M_h(z^{\rm rd})\le\mathcal M_h(z^*/4)\). The two model value comparisons and the \(C_0h\) clipping loss in \(\mathcal M\) give
\[
\mathcal M(z^{\rm cl})\le \mathcal M(z^*/4)+(20C_0^2+C_0)h.
\]
Form \(x^{\rm cl}_i=d_i(1+z_i^{\rm cl})\). Round it up on the physical grid of mesh \(h\), capping at the upper endpoint of the fixed factor-two box in each coordinate. The resulting point is in the fixed box and differs from \(x^{\rm cl}\) by at most \(h\) per coordinate. By the majorant property and the Lipschitz bound on \(F\), the retained physical point has score at most \(\mathcal M(z^*/4)+(20C_0^2+2C_0)h\).

Here is a fixed choice for these bounds. At the start of a homogeneous minimization let \(s\ge4\) bound the binary lengths of the numerators and denominators of its dyadic coefficients and box endpoints, and the binary lengths of its vertex and term counts including the root. Every positive current physical coordinate and every positive given endpoint coefficient lies in \([2^{-s},2^s]\). The relative remainder has weight \(k_e P_i\) and shift \(|P_i-Q_j|/(2P_i)\), where \(P_i=a_e d_i>0\) and \(Q_j=b_e d_j>0\). The weight is at most \(2^{3s}\), the shift at most \(2^{4s}\), and the sum of absolute affine slope contributions of that hinge is at most \(k_e(P_i+Q_j)\le 2^{3s+1}\). An input linear term contributes at most \(2^{2s}\) per coordinate to the sum of absolute slopes. The physical objective has maximum-norm Lipschitz constant at most \(\|\ell\|_1+\sum_e k_e(a_e+b_e)\). These facts and the counts bounded by \(2^s\) give \(C_0=2^{10s}\) as a uniform bound and also give a score range at most \(2^{10s}\). Use the sharper ranges specified in the work count below for iteration counts.

Choose \(h\) as the largest power of two at most \(\min(2^{-s-1},\eta_{\min}/(10^5 C_0^4))\), where \(\eta_{\min}\) is the smallest factor-two target within the invocation. The initial physical endpoints are on this grid, and every retained upper vector is on the same grid. A lower bound of the form \(d_i/2\) is dyadic with at most one more binary digit and is used exactly. Every relative bound and shift is a quotient of dyadics, so integer division gives the desired mesh rounding. If relative clipping selects an exact endpoint, use its corresponding physical endpoint directly. Otherwise \(z_i^{\rm cl}\) is the dyadic from the rounded solve. In both cases the physical-grid operation uses only integers and dyadics of bit length polynomial in \(s+\log(2+1/\eta_{\min})\). For a target \(\eta\) the proved score error is at most \((20C_0^2+2C_0)h<\eta/32\), so the inner gap obeys \(G_{j+1}\le 3G_j/4+\eta/32\). With \(4\lceil\log_2(2+2V/\eta)\rceil\) steps the initial gap contributes at most \(\eta/2\), since \((3/4)^4<1/2\). The sum of the additive errors is at most \(4\eta/32\). Integer comparisons with powers of two give the binary-logarithm ceilings and the dyadic tolerance choices.

\hypertarget{exact-ordinary-calls-and-their-integer-ranges}{%
\subsection*{Exact ordinary calls and their integer ranges}\label{exact-ordinary-calls-and-their-integer-ranges}}

For a directed network of total size at most \(M_0\) with capacities, demands and signed costs integral and bounded in magnitude by a fixed power of \(M_0\), use the deterministic fixed-point algorithm of reference 3, Lemma 2.1. It outputs an explicit integral optimum in \(M_0^{1+o(1)}\) charged word time on feasible instances. Its underlying implementation is Theorem 1.1 of van den Brand, Chen, Kyng, Liu, Peng, Probst Gutenberg, Sachdeva and Sidford. For its ordinary \(\ell_1\) IPM, Theorem 4.6 bounds the operations outside the data structure, the number of data structure queries and the total changed records by \(\widetilde O(m_f\kappa^{-2})\), where \(m_f\) counts arcs and \(\kappa\) is the min-ratio cycle approximation factor. Theorem 5.1 gives \(\kappa=\exp(-O((\log m_f)^{17/18}\log\log m_f))\) and includes initialization and total data structure work, including detection and rebuilding. The flow accuracy prior to rounding is \(1/2\) in cost with exact feasibility. Lemma 4.1 gives a feasible integral flow of no greater cost in \(O(m_f\log(2+m_f))\) work. Applying the rounding conclusion to an optimum also shows that the optimum value is integral, so the rounded cost within \(1/2\) of it is exact. Section 2 gives fixed-point words of \((\log(2+M_0))^{O(1)}\) bits here because the polynomial-range inputs have \(O(M_0\log(2+M_0))\) bits. Reference 3, Lemma 2.1 charges the simulation of these words on logarithmic-bit words. Reference 2, Theorem 6.17, gives the dual-IPM ordinary-flow bound. The finite-word implementation invoked here is the one specified in reference 3. \href{https://arxiv.org/pdf/2407.10830v1}{Reference 2, Theorem 6.17} \href{https://github.com/openai/math/blob/main/preprints/Almost-Linear-Time-Maximum-Cardinality-Matching-in-Sparse-General-Graphs-September-24-2026/main.pdf}{Reference 3, Lemma 2.1} \href{https://arxiv.org/pdf/2309.16629}{Deterministic exact flow, Theorem 1.1, Section 2, Lemma 4.1, Theorems 4.6 and 5.1}

Use the following scaling for the two exact calls. Each prescribed-demand network has a known integral feasible flow \(s^0\). For an arc \(i\to j\) of capacity \(u\), flow \(s^0_e\), and cost \(c_e\), its residual network has a forward arc of capacity \(u-s^0_e\) and cost \(c_e\), and a reverse arc of capacity \(s^0_e\) and cost \(-c_e\). Omit arcs of capacity zero. Adding forward usage and subtracting reverse usage converts any residual circulation into a feasible flow. Conversely, the difference from \(s^0\) of any feasible flow splits into a nonnegative residual circulation using the corresponding direction of each difference. The original cost changes by exactly the residual cost. Thus it suffices to solve a circulation without increasing the capacity or cost magnitudes. We also use this residual construction when augmenting a circulation. An optimal circulation has a residual graph without negative cycles, since a negative directed cycle would permit an augmentation by the minimum of its positive residual capacities. Any nonnegative circulation decomposes into nonnegative multiples of simple directed cycles. To see this, follow positive arcs until a vertex repeats, subtract the smallest flow on the resulting cycle, and repeat on the conserved remaining flow. Each subtraction removes at least one positive arc. This decomposition is used in the analysis only.

For polynomial-range costs and arbitrary integral upper capacities bounded by \(U_f\), start with all capacity digits zero and circulation zero. Reveal the binary capacity digits from most to least significant. At one stage the new capacities are \(2u+\xi\), where the \(\xi_e\) are zero or one and \(u\) is the previous prefix vector. Double the known optimum for \(u\), and form its residual at \(2u+\xi\). Every negative simple residual cycle uses a forward residual arc that was absent before doubling. Indeed, the old optimum had no negative residual cycle, and costs on the old residual arcs are unchanged. Each newly positive arc has capacity one. A minimum-cost residual circulation exists by compactness. Drop all nonnegative cycles from its decomposition to obtain an optimum using only negative simple cycles. The total weight of these cycles is at most the number of newly positive arcs, since each cycle uses at least one such arc. Hence there is an optimum using at most the number of original arcs on every residual arc. Truncate the residual capacities at that number, use Lemma 2.1 with zero demands, and augment by its integral optimum. This yields the new exact optimum since the truncation preserved a comparator of optimal cost. There are \(O(\log(2+U_f))\) stages.

For arbitrary signed integer costs of magnitude at most \(C_f\), keep the full capacities and reveal the bits of the absolute costs, multiplied by the sign of each cost. Start at zero costs with the zero circulation and zero potential. At a stage let \(f\) be the current optimum and \(p\) an integral potential such that residual arcs \(i\to j\) have nonnegative reduced cost \(c^{\rm res}_{ij}+p_i-p_j\). The new cost is twice the old cost plus a signed bit, so with provisional potential \(2p\) the new reduced costs \(r_{ij}\) in the residual of \(f\) are integers at least \(-1\). Let \(v\) be the number of vertices and replace each \(r_{ij}\) by \(\bar r_{ij}=\min(r_{ij},v+1)\). Any simple directed cycle containing an arc with \(r_{ij}>v+1\) now costs at least \(v+1-(v-1)>0\). Solve this bounded-cost residual circulation by the capacity procedure. Its optimum \(\Delta\) has zero flow on every strictly capped arc. Otherwise its cycle decomposition would contain a positive-cost cycle, whose removal would decrease the objective. Thus \(\bar r^\top\Delta=r^\top\Delta\), while \(\bar r^\top g\le r^\top g\) for every feasible residual \(g\). The augmentation is optimal for the true costs as well, since reduced costs and original costs agree on circulation cost.

To renew the potential, take the residual of \(\Delta\) in this capped instance. It has no negative cycle, and the distance-label procedure above gives an integral \(\pi\) using one direct Lemma 2.1 call. Its capacities and demands are at most \(v+1\) and the absolute lengths are at most \(v+1\), so this call has polynomial ranges regardless of \(U_f,C_f\). Set the potential to \(2p+\pi\). For each original arc \(i\to j\) of capacity \(u_e\) and old flow \(f_e\), let \(\Delta_e^+,\Delta_e^-\) be the uses of its forward and reverse residual arcs, extended by zero if absent. Write \(r_e\) for the true forward reduced cost and \(-r_e\) for the reverse. The new forward capacity is \((u_e-f_e-\Delta_e^+)+\Delta_e^-\). If its first term is positive, the capped residual contains the old forward residual arc with length at most \(r_e\). If its second term is positive, the old reverse direction was used and hence uncapped, so its reversal in the capped residual has length \(r_e\). In either case \(\pi_j-\pi_i\le r_e\). The new reverse capacity is \((f_e-\Delta_e^-)+\Delta_e^+\). A positive first term gives the old reverse residual arc of length at most \(-r_e\). A positive second term gives the reversal of a used uncapped forward direction of length \(-r_e\). Either gives \(\pi_i-\pi_j\le-r_e\). These inequalities prove validity for every arc of the true residual. A simple shortest path from the extra source has at most \(v\) arcs, hence \(|\pi_i|\le v(v+1)\). Starting at zero, the recurrence \(p\leftarrow2p+\pi\) over \(O(\log(2+C_f))\) stages keeps bit length \(O(\log(2+M_0)+\log(2+C_f))\). All capacities and maintained flows in the scaling are integer residual differences or prefix values bounded by \(O(U_f)\), and reduced costs have bit length \(O(\log(2+M_0)+\log(2+C_f))\). The general range guarantee is also stated in Appendix C of Chen et al.~\href{https://arxiv.org/pdf/2203.00671v2}{Chen, Kyng, Liu, Peng, Probst Gutenberg and Sachdeva, Appendix C, Lemma C.1}

For an input network with \(M_0\) arcs and vertices in total, each scaling solve uses a graph of size \(O(M_0)\), including any subdivisions before a primitive call. The capacity stages have capacities at most \(M_0\) up to a constant factor, zero demands, and cost magnitude at most the cost stage's vertex count plus one. The distance stages have the polynomial ranges just proved. Take the allowed size in Lemma 2.1 to be a fixed constant times \(M_0+2\) so that it also covers a smaller residual graph with a range bounded using the original size. If an instance has no arcs or all its costs are zero, the known feasible flow suffices. There are \(O(\log(2+U_f)\log(2+C_f))\) bounded-range calls. Each residual and the copy indices for forward, reverse and subdivided arcs are built by a scan. Number its vertices and arcs consecutively for the primitive call. Each augmentation, label traversal and projection takes \(O(M_0)\) arithmetic operations on the stated bit lengths. Arrays indexed by local arc and vertex identifiers and temporary adjacency lists suffice for all this work. This yields \(M_0^{1+o(1)}\operatorname{polylog}(2+U_f+C_f)\) deterministic total work for each exact call, including all recovery.

Each hinge model makes two such exact calls. For the first, choose \(K\) as the smallest power of two strictly above the rounded Lipschitz bound plus one. After scaling by \(1/h\), capacities, net inflow demands and signed costs have magnitude \(O((1+C_0)/h)\), since the root demand is at most the sum of all absolute rounded nonroot demands. For the second, the capacities and demands are at most \(v_o+1\), and the length magnitude is bounded by the first call's cost range \(C_f\). Its scaled distances have magnitude at most \(v_o C_f\). Thus \(U_f,C_f\) can be these bounds rounded up to positive integers. Each model vertex has interval arcs to the root or is the root itself, and each vertex of the label network has an incident source arc. The total sizes are linear in the model vertex and hinge counts even after adding arcs or subdivisions. Building either network and traversing the returned support costs linear work outside the calls, up to the stated arithmetic cost.

\hypertarget{a-feasible-displacement-for-a-bounded-affine-dual}{%
\section*{A feasible displacement for a bounded affine dual}\label{a-feasible-displacement-for-a-bounded-affine-dual}}

Let \(L_e\) be a row with two nonzero entries of opposite signs or one nonzero entry. A unary constraint is one on a single variable. A positive and a negative unary at each vertex bound the feasible dual \(Ly\le c\), whose signed objective is \(b^\top y\). At a strictly feasible \(y\), let \(s_e=c_e-L_ey>0\). The optimum of \(b^\top h\) subject to \([(Lh)_e]_+\le s_e\) is the dual objective gap, since these inequalities are equivalent to feasibility of \(y+h\).

Use a dyadic scale \(\kappa_e\) with \(\kappa_e\le s_e<5\kappa_e\). We adapt the separate signed-value and absolute-motion counters in reference 3, Sections 5.5--5.6, with a single explicit record for each physical dual row. At initialization set the slack anchor \(\hat s_e=s_e\) and both its accumulated signed change \(d_e^{\rm sign}\) and its accumulated absolute change \(d_e^{\rm abs}\) to zero. At every anchoring set \(\kappa_e=2^{\lfloor\log_2\hat s_e\rfloor-1}\). After a dual update \(y\leftarrow y+\Delta y\), add \(-L_e\Delta y\) to the signed change and \(|L_e\Delta y|\) to the absolute change. The exact slack is \(\hat s_e+d_e^{\rm sign}\). If \(d_e^{\rm abs}\ge\kappa_e\), set the anchor to that exact slack, clear both changes and recompute \(\kappa_e\). Between refreshes \(2\kappa_e\le\hat s_e<4\kappa_e\) and \(|d_e^{\rm sign}|\le d_e^{\rm abs}<\kappa_e\), proving the scale bounds. A refresh consumes at least one unit of cumulative absolute slack motion measured in units of that interval's scale, and resetting the counters makes the charged intervals disjoint for each row. A full update takes one scan of rows, with at most one refresh per row. \href{https://github.com/openai/math/blob/main/preprints/Almost-Linear-Time-Maximum-Cardinality-Matching-in-Sparse-General-Graphs-September-24-2026/main.pdf}{Reference 3, Sections 5.5--5.6}

Hold these scales fixed during a displacement search. The normalized matrix has rows \(L_e/\kappa_e\), and the conservative normalized problem uses right hand side one. It contains \(h/5\) for every true feasible displacement \(h\). Define \(h_i^+=[h_i]_+\), \(h_i^-=[-h_i]_+\). For a row \(a h_i-b h_j\) with \(a,b>0\), if \(h_i,h_j\ge0\) or \(h_i,h_j\le0\), one of the row values on \(h^+\) and \(-h^-\) is zero. If \(h_i\ge0\ge h_j\), both row values are nonnegative, and if \(h_i\le0\le h_j\), both are nonpositive. Hence the positive parts add to the positive part on \(h\). This holds for a unary because one value is zero. Both pieces of a conservative displacement are therefore conservative. Their signed objective scores sum to its score. From an optimal true displacement this proves that one of two nonnegative-vector problems has optimal score at least a tenth of the true objective gap. In the negative problem use \(x=h^-\), with every row and the score negated. Drop the nonpositive scalar constraints in each problem. Write each remaining pair as \(a_e x_i-b_e x_j\le1\) and each upper unary as \(a_e x_i\le1\), with positive coefficients on their actual endpoints. There is an upper unary at every vertex, so each problem has an attained maximum on a compact set.

\hypertarget{a-coordinatewise-upper-point-with-bounded-violation}{%
\subsection*{A coordinatewise upper point with bounded violation}\label{a-coordinatewise-upper-point-with-bounded-violation}}

Let \(n_0\) count the coordinates and \(m_0\ge n_0\) bound the number of retained rows. In the application \(m_0\) is the total augmented row count before discarding the nonpositive unaries of either sign problem. For \(n_0=0\) the empty displacement suffices. Otherwise let \(H\ge16\) be a power of two bounding \(m_0,n_0\), the positive dyadic coefficients of the retained rows, their reciprocals and their dyadic denominators. Take the smallest power-of-two integer \(p\ge p_0=\max(4,\lceil\log_2(2m_0)\rceil)\). Then \(p<2p_0\), and \(p\le H\) since \(H\) is a power of two with \(H\ge p_0\). Define \(E(x)=\sum_e[(L_{\rm norm}x)_e]_+^p\), where the sum uses the retained normalized rows in this sign problem. With \(v_i=x_i^p\) define \(\mathcal E(v)=E((v_i^{1/p})_i)\) on \(v\ge0\). A pair contributes
\[
 \mathcal E_e(v)=\big((a_e^p v_i)^{1/p}-(b_e^p v_j)^{1/p}\big)_+^p .
\]
For \(Y=b_e^p v_j>0\) and \(X=a_e^p v_i\ge0\), this is \(Y\phi(X/Y)\), where \(\phi(z)=(z^{1/p}-1)_+^p\). The derivative is zero up to 1 and then \((1-z^{-1/p})^{p-1}\), which increases, giving convexity of \(\phi\). In a convex combination of two pairs \((X_1,Y_1),(X_2,Y_2)\) with positive \(Y_1,Y_2\) and weights \(\lambda,1-\lambda\), the combined ratio \(X/Y\) averages \(X_1/Y_1,X_2/Y_2\) with weights \(\lambda Y_1/Y,(1-\lambda)Y_2/Y\). Convexity of \(\phi\) and multiplication by \(Y\) give convexity of the pair contribution. The original pair formula is continuous at \(Y=0\) with value \(X\), so the inequality extends there by taking limits. A positive unary contributes \(a_e^p v_i\).

We compute a convex homogeneous hinge sum \(Q\) such that \(Q\le\mathcal E\le4Q\) on \(v\ge0\). Index breakpoints by \(k\) and use \(z_k=(1+\Delta_k)^p\) with \(\Delta_0=H^{-12}\) and \(\Delta_{k+1}=(1+1/p)\Delta_k\), through the first \(\Delta_J\ge10p\). The unrounded slope at \(z_k\) is \(t_k=(\Delta_k/(1+\Delta_k))^{p-1}\). The staircase has slope zero below \(z_0\) and slope \(\tilde t_k\) on each interval starting at \(z_k\), where \(\tilde t_k\) floors \(t_k\) to mesh \(H^{-10p}\). For \(k<J\), \(t_{k+1}/t_k\le(1+1/p)^{p-1}<3\), so \(\phi'\le3\tilde t_k+3H^{-10p}\) there. On the final interval \(t_J\ge(1-1/(1+10p))^{p-1}\ge1/2\), and \(\phi'\le1\). Below \(z_0\), \(\phi'\le H^{-12(p-1)}\le H^{-6p}\). Thus the derivative of the integrated staircase \(\tilde\phi\) satisfies \(0\le\tilde\phi'\le\phi'\le3\tilde\phi'+3H^{-6p}\) off the breakpoints. Integrating from zero gives the claimed scalar lower approximation and \(\phi(z)\le3\tilde\phi(z)+3H^{-6p}z\).

For each pair \(e=(i,j)\), taking the perspective of \(\tilde\phi\) gives the hinges \([a_e^p v_i-z_k b_e^p v_j]_+\) for \(0\le k\le J\), with weights \(\tilde t_k-\tilde t_{k-1}\ge0\) and \(\tilde t_{-1}=0\). Retain the unary terms exactly in \(Q\). The total extra linear error on pairs is at most \(3H^{-6p}\sum_{e=(i,j)} a_e^p v_i\le3m_0H^{-4p}\sum_i h_i^p v_i\), where \(h_i\ge H^{-1}\) is a chosen upper unary coefficient at \(i\). Since \(3m_0H^{-4p}<1\), adding all unary contributions proves \(\mathcal E\le4Q\). The inequalities include zero endpoints by continuity of the pair formula and hinges. The breakpoint count is \(O(p\log(Hp))\), using \(\log(1+1/p)\ge1/(2p)\). If \(p=2^r\), then \(\Delta_k=H^{-12}(p+1)^k/2^{rk}\), so forming each dyadic \(z_k\) uses binary length \(O(p(\log H+(k+1)\log(p+1)))\). Integer powers, a division to floor each slope, and differences of consecutive floored slopes give the model. Store only positive increments. Each computation uses polynomially many bits in \(p,\log H\).

On \([H^{-6p},H^{6p}]^{n_0}\) minimize \(J(v)=Q(v)+\sum_i B(v_i)\) to error \(1/64\), with \(B(1)=0\) and slope \(-2^{-j}\) on \([2^j,2^{j+1}]\). Put \(J_B=6p\log_2 H\). The initial slope is \(-2^{J_B}\), and each interior breakpoint \(2^j\), \(-J_B<j<J_B\), increases the slope by \(2^{-j}\). Thus this convex piecewise-linear regularizer is implemented by the initial affine term and one rooted hinge per interior breakpoint. At \(t=2^j r\) for \(1\le r\le2\) it has value \(-j-(r-1)\), since each full dyadic interval decreases its value by one. It is decreasing, \(|B(t)|\le J_B\) on the box, and \(B(2t)=B(t)-1\) when both arguments are in the box. Comparing to \(\mathbf1\) gives \(J(v)\le m_0H^p+1/64\). Because \(Q(v)\ge H^{-p}\sum_i v_i\) by the upper unaries, we have
\[
 H^{-p}v_i\le m_0H^p+6n_0p\log_2H+1/64
 \le H^{p+1}+6H^3+1/64<H^{5p}/4 .
 \]
The last inequality uses \(H\ge16,\ p\ge4\). If \(v_i\le2H^{-6p}\), doubling that coordinate stays in the box and saves one in \(B\). By homogeneity and subadditivity of \(Q\), its increase is at most \(Q(v_i\mathbf e_i)\le m_0H^p v_i\). Subadditivity follows termwise from \([a+b]_+\le[a]_++[b]_+\). The increase is at most \(2m_0H^{-5p}<1/2\), contradicting the \(1/64\) optimality error. Thus \(v/2\) is in the box. Its comparison, \(J(v)\le J(v/2)+1/64\), gives \(Q(v)/2\le n_0+1/64\). For any nonnegative \(x\) feasible for the normalized right hand sides one, let \(w_i=x_i^p\). Then \(Q(w)\le E(x)\le m_0\) and each \(w_i\le H^p\) by its upper unary. The trial \(v+w/(16m_0)\) is in the box since \(v_i<H^{6p}/4\) and \(w_i\le H^p\). It increases \(Q\) by at most \(1/16\). If it at least doubled any coordinate, it would decrease \(\sum_i B(v_i)\) by at least one. Thus each \(x_i^p\le16m_0 v_i\).

Round the \(p\)-th roots of \(v_i\) upward to absolute error at most \(1/(100m_0H)\) and divide by 64. Use a dyadic mesh no coarser than this error. Find each ceiling by integer bisection on the grid in \([0,H^6]\), comparing the exact \(p\)-th power of the tested point with \(v_i\). Rounding increases any positive row value by at most \(1/(100m_0)\), since its positive endpoint has coefficient at most \(H\) and increasing its negative endpoint can only help. The triangle inequality in \(\ell_p\) and \(\mathcal E(v)\le4Q(v)\) bound the positive-row norm before division by \((8n_0+1/8)^{1/p}+m_0^{1/p}/(100m_0)<16\), using \(m_0^{1/p}\le2,\ p\ge4\). Also \((16m_0)^{1/p}\le8\). Thus the resulting \(D\) is feasible for the normalized right hand sides one, and \(D_i\ge x_i/A\) for every such nonnegative feasible \(x\) and all \(i\), with \(A=512\). In particular its coordinates obey their actual normalized upper unary limits. Its construction gives an upper box even for signed linear objectives.

\hypertarget{propagating-the-affine-allowance-through-decreasing-upper-bounds}{%
\subsection*{Propagating the affine allowance through decreasing upper bounds}\label{propagating-the-affine-allowance-through-decreasing-upper-bounds}}

For the nonnegative problem let \(g^\top x\) be its signed score, where \(g=b\) for positive displacements and \(g=-b\) for negative displacements. Specify the number \(T\) of box steps and set \(\epsilon\) to the largest power of two at most \(1/(8T)\), so \(1/(16T)<\epsilon\le1/(8T)\). Put \(\varepsilon=\min(\epsilon,1/(2A))\). Choose a positive floor \(l\le\min_i D_i\) with \(a_e l\le\varepsilon\) for every row. Start at \(d=D,\tau=1\). For analysis set the shadow point \(u=l\mathbf1\vee(\varepsilon x^*)\), where \(x^*\) maximizes the conservative nonnegative score with right hand side one. Both vectors in that maximum satisfy the row allowance \(\varepsilon\). In any pair row the maximum has the source value of one of the two vectors and a destination value at least that vector's destination value, so their maximum also satisfies the allowance. For a unary use its source alone. Since \(D_i\ge x_i^*/A\), the shadow obeys \(l\le u\le d\). Hence its initial score is at least \(\varepsilon g^\top x^*-\|g\|_1l\), because the move from \(\varepsilon x^*\) to \(u\) has maximum-norm length at most \(l\). Maintain at each row
\[
 a_e d_i-b_e d_j\le\tau,\qquad
 a_e u_i-b_e u_j\le\max(\varepsilon,\tau u_i/d_i).                    \tag{1}
\]
An absent destination contributes zero in all formulas.

In the box \([\max(l,d/2),d]\) form the homogeneous inequalities
\[
                   (a_e-(\tau+\epsilon)/d_i)_+ z_i\le b_e z_j .       \tag{2}
\]
For nonnegative \(z\), (2) is equivalent to \(a_e z_i-b_e z_j\le(\tau+\epsilon)z_i/d_i\). If \(a_e-(\tau+\epsilon)/d_i\le0\), both forms hold from nonnegativity of the destination contribution, including for a unary. The upper point \(d\) satisfies the latter inequality by (1). For the midpoint \(w=(d+u)/2\), (1) gives \(a_e w_i-b_e w_j\le\tau w_i/d_i+\varepsilon/2\), using \(\max(\varepsilon,\tau u_i/d_i)\le\varepsilon+\tau u_i/d_i\). Since \(w_i/d_i\ge1/2\) and \(\varepsilon\le\epsilon\), the midpoint satisfies (2). Floor the left coefficients of (2) to nonnegative dyadics \(\bar a_e\) with their errors times \(D_{\max}=\max_i D_i\) at most \(\epsilon/16\). Maximize
\[
         g^\top z-M_b\sum_e[\bar a_e z_i-b_e z_j]_+
\]
on the box to additive \(\eta\), retaining a feasible box point \(v\). Choose a dyadic \(M_b>16(1+2n_0\|g\|_\infty D_{\max}+\eta)/\epsilon\) within a factor two of this threshold. Apply the homogeneous hinge minimizer to the negative of this objective. The penalty at the midpoint is zero, giving \(g^\top v\ge(g^\top d+g^\top u)/2-\eta\) because the penalty at \(v\) is nonnegative. Comparison with \(d\), where the penalty is zero, bounds \(M_b\) times the sum of rounded violations by \(g^\top(v-d)+\eta\le2n_0\|g\|_\infty D_{\max}+\eta\). Thus that sum is at most \(\epsilon/16\).

The point \(v\vee u\) is in the current box and has each rounded violation at most that of \(v\). If the source comes from \(v\), the destination contribution can only increase. Otherwise \(u_i\ge v_i\ge d_i/2\), so \(\tau u_i/d_i\ge1/2>\varepsilon\) and (1) ensures \(u\) satisfies (2) for that row. The maximum retains this source and can only increase the destination contribution. Write \(R(z)=\sum_e[\bar a_e z_i-b_e z_j]_+\). Approximate optimality gives \(g^\top v\ge g^\top(v\vee u)+M_b(R(v)-R(v\vee u))-\eta\ge g^\top(v\vee u)-\eta\). Since \((v\vee u)+(v\wedge u)=v+u\) coordinatewise, \(g^\top(v\wedge u)\ge g^\top u-\eta\).

Update \(d'=v,\ u'=u\wedge v,\ \tau'=\tau+2\epsilon\). The violations and coefficient rounding give
\[
 a_e v_i-b_e v_j\le(\tau+\epsilon)v_i/d_i+\epsilon/8\le \tau'v_i/d_i
\]
because \(v_i/d_i\ge1/2\). Thus the physical row of \(d'\) is at most \(\tau'\). Both \(v\) and \(u\) satisfy \(a_e x_i-b_e x_j\le\max(\varepsilon,\tau' x_i/d_i)\). Equivalently, the destination contribution is at least \(\min(a_e x_i-\varepsilon,(a_e-\tau'/d_i)x_i)\). If the second slope is negative, the right hand side of this last inequality is nonpositive and the condition holds for all nonnegative \(x\). Otherwise the right hand side is nondecreasing in \(x_i\). For a coordinatewise minimum take the vector attaining its smaller destination. Lowering that vector's source to the minimum of the two source coordinates preserves the condition. For a unary the destination contribution is always zero and the same monotonicity proves the condition at the smaller source. Hence \(u'\) satisfies the condition, and replacing \(d_i\) by \(v_i\le d_i\) in the denominator in (1) relaxes it. This proves (1) inductively.

After \(T\) steps \(\tau\le5/4\). With \(q_0=g^\top u_0\), the shadow score at step \(j\) is at least \(q_0-j\eta\) by induction. For \(G_j=q_0-g^\top d_j\), where \(d_j\) here denotes the vector at step \(j\), the midpoint comparison gives \(G_{j+1}\le G_j/2+(j/2+1)\eta\). Initially \(G_0\le\|g\|_1D_{\max}\). Induction gives \(G_T\le2^{-T}\|g\|_1D_{\max}+(T+1)\eta\), and thus
\[
 g^\top d_T\ge \varepsilon g^\top x^*-\|g\|_1 l
       -2^{-T}\|g\|_1D_{\max}-(T+1)\eta.                             \tag{3}
\]
Thus uniform positive affine row allowances suffice with a single inverse-logarithmic loss. All computations concern the original physical directions.

\hypertarget{a-bounded-additive-solver-with-primal-output}{%
\section*{A bounded additive solver with primal output}\label{a-bounded-additive-solver-with-primal-output}}

The input here has dyadic positive widths \(U_a\), nonnegative dyadic costs \(q_a\), an attainable dyadic right hand side \(d_{\rm box}\), and dyadic matrix entries of magnitude at most one. Each nonzero column has either two entries of opposite signs or one signed entry. Write \(\mathrm{OPT}_{\rm box}\) for the minimum cost under exact equations and box constraints. A free zero-cost signed slack in \([-s,s]\) is represented by the difference of two nonnegative variables in \([0,s]\) on opposite unit columns. Every such difference is in \([-s,s]\), and any prescribed member of \([-s,s]\) is obtained by taking its positive and negative parts. Set a variable on a zero column or with zero width to zero. This keeps the equation vector unchanged and can only lower the nonnegative objective, so the resulting optimum is the same. If there are no equations, all variables zero gives cost zero. Let \(0<\epsilon_{\rm out}\le1\) be the dyadic target for the maximum equation error and the additive objective error, and set \(C=1+\sum_a q_a U_a\). The solver returns exact box coordinates with both errors at most \(\epsilon_{\rm out}\).

Use the capacity conversion of reference 2, Appendix B, Definition B.2 and Lemma B.5, with the actual endpoint coefficients. For a column \(\alpha_a\mathbf e_{i(a)}-\beta_a\mathbf e_{j(a)}\), add a private vertex \(k(a)\), and use unbounded nonnegative variables \(x_a,z_a\) on \(\alpha_a\mathbf e_{i(a)}-\mathbf e_{k(a)}\) and \(\beta_a\mathbf e_{j(a)}-\mathbf e_{k(a)}\), with costs \(q_a,0\). A missing endpoint of a unary has zero contribution in these formulas. In particular a positive unary gives a column for \(x_a\) with its positive entry and the private negative entry, and a private-only column for \(z_a\). For a negative unary it is \(x_a\) that has only the private entry. Start with \(d_{\rm box}\) at the old vertices, add \(\beta_a U_a\) at \(j(a)\) when present, and set the private demand at \(k(a)\) to \(-U_a\). The private equality now gives \(x_a+z_a=U_a\), and the contributions at the old vertices are the original column times \(x_a\) plus \(\beta_a U_a\) at \(j(a)\). This proves equivalence before artificial columns, in both directions, at the same cost. \href{https://arxiv.org/pdf/2407.10830}{Reference 2, Appendix B}

Add positive and negative unit columns at every augmented vertex, at cost \(H_a\) for each, where \(H_a\) is the smallest power of two at least \(64(C+1)/\epsilon_{\rm out}\). Denote the resulting equations by \(L^\top f=b,\ f\ge0\) and the cost vector by \(c\), with \(m_0\) rows of \(L\) and \(n_0\) augmented vertices. Each exact box solution extends at its cost, with \(z_a=U_a-x_a\), expanded mass \(\sum_a U_a\), and artificial mass zero. The dual is to maximize \(b^\top y\) subject to \(Ly\le c\). For its strict seed \(y^{\rm seed}\), set old potentials to zero and private potentials to one. Every nonartificial row evaluates to \(-1\), and every artificial row to zero or \(\pm1\), with \(H_a\ge2\). Thus every slack is at least one and the artificial inequalities bound all feasible dual coordinates.

In scans of these unperturbed data form the counts, total widths, \(C,H_a,\|c\|_\infty,\|b\|_\infty\). Use integer bit lengths and comparisons with powers of two to choose the smallest power of two \(P\ge2^{20}\) bounding these counts and total widths, \(C+2,H_a+2,\|c\|_\infty+2,\|b\|_\infty+2,\epsilon_{\rm out}^{-1}\) and the dyadic denominators and inverse nonzero magnitudes of the unperturbed data. In the application also require \(P\ge N\). This is polynomial in the normalized input size parameter used below. For each trial set \(\sigma=P^{-12}\), sample independent uniform increments on the mesh \(P^{-160}\) in \([\sigma,2\sigma)\) for costs and \([0,\sigma)\) for demands, and write \(c',b'\) for the perturbed vectors. Each sample uses exactly \(148\log_2 P\) independent random bits, interpreting the resulting integer as a multiple of the mesh and adding \(\sigma\) for costs. Let \(\mathrm{OPT}'=\min\{(c')^\top f:L^\top f=b',\ f\ge0\}\). Add to an exact box comparator the artificial unit flows that supply all demand increments. Therefore
\[
 \mathrm{OPT}'\le \mathrm{OPT}_{\rm box}
                      +2\sigma\sum_a U_a+n_0(H_a+2\sigma)\sigma .     \tag{4}
\]
The perturbed primal is feasible by this construction and attains its minimum since \(c'_e\ge\sigma>0\) and a finite cost sublevel in the nonnegative orthant is compact. In (4), \(\sum_a U_a\le P,\ n_0\le P\), and \(H_a+2\sigma\le P\), so the increment is at most \(2P^{-11}+P^{-10}<1\). Since \(\mathrm{OPT}_{\rm box}\le C-1\le P-3\), we have \(\mathrm{OPT}'<P\), and every optimal coordinate is at most \(P/\sigma\le M=P^{15}\). The perturbed artificial inequalities bound every dual coordinate by \(2P\), and the original seed still has slacks at least one. The perturbations and the margins in the choice of \(P\) give \(\|c'\|_\infty\le P,\ \|b'\|_\infty\le P,\ \|b'\|_1\le P^2\). Strong LP duality identifies \(\mathrm{OPT}'\) with the attained maximum of \((b')^\top y\).

Initialize \(y=y^{\rm seed}\) and set the anchors from \(s=c'-Ly^{\rm seed}\). At each iterate use both displacement searches with scores \(g=\pm b'\) and with \(H=P^{512},\ T=10^4\log_2 P,\ \eta=P^{-150},\ l=P^{-10000}\) in (1)--(3). Choose the better signed endpoint or zero by its exact score, and move by one quarter of that signed displacement. At the final upper point (1) bounds each normalized row value by \(5/4\) on the searched sign. Rows dropped for that sign have nonpositive value. Therefore each physical row increase is at most \(5\kappa_e/16\le5s_e/16\), leaving slack at least \(11s_e/16\). Mix this potential toward \(y^{\rm seed}\) using weight \(\nu=P^{-200}\) on the seed, then round coordinates to the nearest point on mesh \(P^{-210}\). The rounding costs at most \(2P^{-210}\) in slack since every row has at most two coefficients of magnitude at most one. Thus each new slack is at least
\[
(1-\nu)(11/16)s_e+\nu-2P^{-210}\ \ge\ (2/3)s_e+\nu/2,
\]
using \(\nu\le 1/33\) and \(2P^{-210}\le\nu/2\). In particular every stored iterate has slacks at least \(P^{-200}/2\), remains dual feasible, and has slacks at most \(5P\) by the artificial unary bounds and \(\|c'\|_\infty\le P\).

Update the exact slack anchors after this rounding. The row coefficients are multiples of \(1/P\), so each change is a multiple of \(P^{-211}\), and the perturbed cost grid is a multiple of that mesh as well. The absolute change in one row between stored iterates is at most \(8P\) because the two iterates both lie in \([-2P,2P]^{n_0}\). Thus the accumulators and anchors have \(O(\log P)\) bits over the \(O(\log^2 P)\) scans established below. The dyadic scales also approximate the dual \(\ell_1\) lengths \(s_e^{-1-\alpha}\) from the slack scaling of reference 2, Theorem 6.3. For \(\alpha=1/(500\log_2 P)\), the interval \([P^{-201},P^2]\) for \(s_e\) gives \(1/2\le s_e^\alpha\le2\). Consequently \(\kappa_e^{-1}/s_e^{-1-\alpha}=(s_e/\kappa_e)s_e^\alpha\) lies in \([1/2,10]\). The displacement oracle uses these dyadic inverse scales in its rowwise allowances. \href{https://arxiv.org/pdf/2407.10830v1}{Reference 2, Theorem 6.3 and Section 6.1}

The bounds \(\kappa_e\le s_e<5\kappa_e\) give \(1/(5P)\le1/\kappa_e<10P^{200}\). The original nonzero coefficient magnitudes belong to \([1/P,1]\), with dyadic denominators at most \(P\). Multiplying by the power of two \(1/\kappa_e\) gives a normalized magnitude in \([1/(5P^2),10P^{200}]\). Its denominator increases by at most \(5P\) when \(\kappa_e>1\) and can only decrease when \(\kappa_e\le1\), so the denominator is at most \(5P^2<H\). Use \(P^{-12000}\) for the center's root rounding, which is at most \(1/(100m_0H)\) since \(m_0\le P\). The center gives \(D_i\ge H^{-6}/64>l\). The retained artificial upper unary at a vertex has coefficient \(1/\kappa_e\), so its feasibility inequality gives \(D_i\le\kappa_e\le s_e\le H_a+2\sigma+\|y\|_\infty\le4P\) on each sign.

Use left-coefficient mesh \(P^{-32}\) in (2), rounding the nonnegative part downward by dyadic quotient and integer floor. For \(P=2^k\) with \(k\ge20\), \(10^4k\le2^k\) holds at \(k=20\) and remains true when \(k\) increases by one, so \(T\le P\). We have \(1/\varepsilon\le16T\) by \(1/\epsilon<16T\) and \(2A=1024\le16T\). Also \(M_b\le32(2+8P^3)16P<P^6\) by \(\|g\|_\infty\le P,\ n_0\le P,\ D_{\max}\le4P\), and \(1/\epsilon<16P\). The coefficient rounding changes each row by at most \(4P\cdot P^{-32}<1/(256P)<\epsilon/16\), and \(Hl=P^{-9488}<1/(16P)\le\varepsilon\). To keep the box endpoints on a fixed grid across the \(T\) steps, solve each box to error \(\eta/2\) and round its returned physical point upward to the \(P^{-30000}\) grid, capped at the old upper point. This grid contains the floor and the initial \(D_i\), including the division by 64. The score change is at most \((P^2+2P\cdot P^6H)P^{-30000}<\eta/2\), using the maximum-norm Lipschitz bound \(\|g\|_1+2m_0M_bH\). Thus the retained \(v\) has the accuracy used in (1)--(3). All coefficients, endpoints and tolerance bounds used to form models are dyadic, with the inner model shifts rounded by the earlier fixed-mesh procedure.

Let \(G=\mathrm{OPT}'-(b')^\top y\). One sign has conservative optimum score at least \(G/10\). By (3) and the quarter step, the objective improves by at least \(\varepsilon G/40\) minus the losses in (3), bounded by \(P^2 l+4P^3 2^{-T}+(T+1)P^{-150}\). Here \(l=2^{-T}=P^{-10000}\), so these losses total at most \(5TP^{-150}\). Taking the best of both signs and zero preserves this lower bound. The seed mixture changes the score by at most \(\nu\|b'\|_1(2P+1)\le3P^{-197}\), and rounding costs at most \(P^{-210}\|b'\|_1\le P^{-208}\). Thus, with \(L_{\rm err}=5TP^{-150}+16P^{-197}\), each full update gives \(G_{\rm new}\le(1-\varepsilon/64)G+L_{\rm err}\). Initially \(G_{\rm seed}\le P+\|b'\|_1\le8P^3\). With \(R=\lceil 20000\log_2 P/\varepsilon\rceil\) rounds, the remaining gap is at most
\[
 8P^3 \exp(-20000\log_2 P/64)+(64/\varepsilon)L_{\rm err}<P^{-120}.
\]
In the first term \(\log_2 P\ge\log P\), with \(\log\) the natural logarithm in this comparison. The first term is at most \(8P^{3-20000/64}<P^{-120}/2\). In the second, \(64/\varepsilon\le1024T\le1024P\) bounds the term by \(5120P^{-148}+16384P^{-196}<P^{-120}/2\). Write \(\delta_0=P^{-120}\) for this dual accuracy. The choices of \(T,\varepsilon\) give \(R=O(\log^2 P)\).

The total number of scale changes is at most \(m_0R\), since there is at most one change per row in a round. The regularization also bounds absolute logarithmic slack motion. Each decrease in slack is by a factor of at most \(3/2\), so its cumulative downward logarithmic motion is at most \(R\log(3/2)\) per row. The upper and lower slack limits bound the net logarithmic change in absolute value by \(\log(10P^{201})\). The cumulative upward motion is therefore at most \(R\log(3/2)+\log(10P^{201})\). On an interval that triggers a refresh, the slack before the final update is below \(5\kappa_e\). If its final value exceeds \(6\kappa_e\), the last update spends logarithmic motion greater than \(\log(6/5)\). If the final value is at most \(6\kappa_e\), the whole interval lies below \(6\kappa_e\). For two slacks in that interval, \(|\log s'-\log s|\ge |s'-s|/(6\kappa_e)\) by integrating \(1/s\). The refresh trigger accumulated at least \(\kappa_e\) in absolute motion, so in this case the interval spends at least \(1/6\) in logarithmic motion. Each refresh thus charges a constant amount on disjoint intervals. The \(O(m_0(R+\log P))\) motion-based recourse bound and the \(O(m_0 R)\) explicit scan count both apply to sparse graphs.

\hypertarget{extracting-an-approximately-feasible-primal-on-independent-columns}{%
\subsection*{Extracting an approximately feasible primal on independent columns}\label{extracting-an-approximately-feasible-primal-on-independent-columns}}

For analysis, keep the box \([-2P,2P]^{n_0}\), redundant in the full dual, and delete the inequality \(e\) from \(Ly\le c'\). Call the resulting compact nonempty polytope \(K_{-e}\). For a fixed \(e,i\), condition on all costs and all objective coordinates except \(i\), and let \(\psi(t)\) be the optimal value on \(K_{-e}\) at \(b'_i=t\). It is convex as a maximum of affine functions and \(2P\)-Lipschitz because every candidate slope satisfies \(|y_i|\le2P\). For \(\omega=P^{-65}\), evaluating any \(\delta_0\)-optimal point \(z\) at \(t-\omega\) and \(t+\omega\) gives
\[
  \frac{\psi(t)-\psi(t-\omega)-\delta_0}{\omega}
  \le z_i\le
  \frac{\psi(t+\omega)-\psi(t)+\delta_0}{\omega}.
\]
Put \(q=P^{-160}\) and \(t_j=b_i+jq\), so the uniform sample has \(0\le j<\sigma/q\). Write \(A_j=(\psi(t_j+\omega)-\psi(t_j))/\omega\), extended to all integer \(j\). The backward quotient at \(t_j\) is \(A_{j-\omega/q}\). When summed over the sampling interval, these two sequences cancel except for \(\omega/q\) terms at each end. Each \(|A_j|\le2P\), and the difference of quotients is nonnegative by convexity. Consequently the expected length of the coordinate interval in the display is at most \(4P\omega/\sigma+2\delta_0/\omega\). Markov's inequality and a union bound over \(m_0n_0\le P^2\) pairs show that with failure probability at most \(P^2(4P\omega/\sigma+2\delta_0/\omega)/\rho\le 6P^{-20}\), all \(\delta_0\)-optimal points of each \(K_{-e}\) differ by at most \(\rho=P^{-30}\) in each coordinate.

Read \(S=\{e:c'_e-L_e y\le\theta\}\) with \(\theta=P^{-24}\). Call \(e\) dispensable if some optimal primal \(f^*\) has \(f^*_e=0\). For such a witness and every \(z\in K_{-e}\), \((b')^\top z=\sum_{a\ne e} f_a^*L_a z\le(c')^\top f^*\). Thus deleting \(e\) preserves the dual optimum and \(y\) is \(\delta_0\)-optimal on \(K_{-e}\). For each \(e\), take the lexicographically least optimizer \(y^{-e}\) there as an analysis witness, obtained by minimizing successive coordinates on the compact set of optimizers. Conditional on the remaining costs and demands, \(y^{-e}\) is fixed independently of \(c'_e\). On the width event a dispensable \(e\in S\) therefore implies \(|c'_e-L_e y^{-e}|\le\theta+2\rho\), using \(\|L_e\|_1\le2\). An interval of length \(2(\theta+2\rho)\) contains at most \(2(\theta+2\rho)/q+2\) points of the cost grid, which has \(\sigma/q\) samples. Bound this interval event under the remaining costs and demands alone. Summing its probability and adding the width failure gives total probability at most
\[
6P^{-20}+P\big(2(\theta+2\rho)/\sigma+2q/\sigma\big)
=6P^{-20}+2P^{-11}+4P^{-17}+2P^{-147}<10P^{-11}.
\]
On the complementary event, called the good event below, every row of \(S\) is indispensable.

On the resulting good event all columns in \(S\) are independent. Indeed, fix any exact optimal primal \(f^*\) and dual \(y^*\). Then \(f^*_e>0\) and \(c'_e=L_e y^*\) on \(S\). If a nonzero \(w\) supported on \(S\) had \(L^\top w=0\), choose its sign to have a positive entry and set \(t=\min_{w_e>0}f_e^*/w_e>0\). The vector \(f^*-tw\) is nonnegative, satisfies all primal equations, and has the same cost because \((c')^\top w=(y^*)^\top L^\top w=0\). It vanishes on at least one column of \(S\), a contradiction. Also
\[
 \sum_e f_e^*(c'_e-L_e y)=\mathrm{OPT}'-(b')^\top y\le\delta_0 .
\]
Every term is nonnegative, and outside \(S\) its slack factor exceeds \(\theta\). The mass there is at most \(\delta_0/\theta=P^{-96}\). Keep \(f^*\) on \(S\) and zero elsewhere to obtain a witness in \([0,M]\) of maximum vertex residual at most \(P^{-96}\), since the coefficients have magnitude at most one.

Build the undirected multigraph on the augmented vertices using two-entry columns of \(S\) as edges, while recording unary columns at their single vertices. There are no loops because each two-entry column has distinct endpoints. Each connected component on \(q_v\) vertices has at most \(q_v\) retained columns by independence, as all their entries are supported on those vertices. A spanning tree uses \(q_v-1\) pairs. Thus there are either exactly \(q_v-1\) pairs and at most one unary or \(q_v\) pairs and no unary. Include \(q_v=1\). Test these counts during adjacency-list traversal before running the recovery. A failed test rejects the trial. Model a unary in a tree as an edge to a private leaf with no demand equation, and root each tree at a constrained vertex.

On a tree and trial tolerance \(d_0\), sweep from leaves to root. A nonroot message is an interval of values of the parent edge. A dummy leaf sends \([0,M]\). At each constrained vertex \(i\), let \([L_c,U_c]\) be the sum of the images of the already computed child intervals under their signed coefficients at \(i\). The child intervals are independent, and the possible sums fill the summed interval. Use \([0,0]\) for no children. With demand \(d_i^{\rm tree}\) and nonzero signed coefficient \(\alpha\) on the parent edge, compute the exact interval \(I\) of parent flows in \([0,M]\) with
\[
 \alpha x_{\rm parent}\in[d_i^{\rm tree}-d_0-U_c,\ d_i^{\rm tree}+d_0-L_c].
\]
This uses division by \(\alpha\) with the endpoints ordered. Reject if the interval is empty. Otherwise send its outward rounding on mesh \(h_{\rm tree}=P^{-120}\), clipped to \([0,M]\), and retain \(I\) for the backward pass. At the root check that \([L_c,U_c]\) meets \([d_i^{\rm tree}-d_0,d_i^{\rm tree}+d_0]\), which also handles an isolated root with zero child sum. By induction an exact feasible choice for tolerance \(d_0\) survives every forward test, because the stored child intervals then contain its child flows.

On acceptance perform a root-to-leaf pass. A stored parent value lies in the rounded message, so it is within \(h_{\rm tree}\) of its projection to \(I\). At that projected value, choose a child contribution sum in \([L_c,U_c]\) that meets the local tolerance after the projected parent contribution. The root does this directly without a parent. Start each child at its smallest contribution and distribute the needed excess up to the maximum contribution for that child, using whole endpoints until at most one partial child value. Divide by that child's signed coefficient to get its flow, and round it inside its stored interval to the mesh. The mesh contains each stored interval endpoint and \(M\), so this costs at most \(h_{\rm tree}\) in flow. Restore the actual assigned parent value. Because all coefficient magnitudes are at most one, the projection and flow rounding together incur at most \((\deg(i)+1)h_{\rm tree}\) in error at \(i\). Each selected child value is stored exactly and used as the actual parent value in its child's pass. Hence the error bound is local and incurs no multiplicative factors along paths. All sums in the forward pass use dyadic child messages, each exact message endpoint uses a single division, and only grid values are sent. In the backward pass the exact projection and at most one partial child division at each vertex likewise use bit lengths polynomial in \(\log P\) on the demand and tolerance grids specified below.

In a unicyclic component remove the unique edge outside the spanning tree and fix its flow to \(t\in[0,M]\), subtracting its column contributions from the two endpoint demands. This applies to parallel edges as well. Write \(F(t)\) for the minimum maximum residual on the remaining tree over its box of flows. This minimum exists by compactness. Averaging optimal tree flows for two parameters proves convexity, and keeping the same tree flow while varying \(t\) changes its maximum residual by at most the parameter change. Thus \(F\) is 1-Lipschitz. For fixed \(t\), start with a bracket \([0,4P^{17}]\) for \(F(t)\), which contains the residual of the all-zero tree flow. Use the exact dyadic midpoint \(d\) as the tree tolerance. If rejected, replace the bracket's lower endpoint by \(d\), since an exactly \(d\)-feasible choice survives. Otherwise cap the upper endpoint at \(d+(P+1)h_{\rm tree}\) by the backward error bound. Each bracket width is at most half the previous width plus \((P+1)h_{\rm tree}\). After \(100\log_2 P+5\) steps its width is bounded by \(4P^{17}2^{-(100\log_2P+5)}+2(P+1)h_{\rm tree}<P^{-80}\). Return the upper endpoint as an approximation with error at most \(\mu=P^{-80}\). The demands and all these trial tolerances have \(O(\log P)\) bits since \(t\) is on the grid specified next, and each bracket update uses one dyadic halving and an addition of a fixed mesh multiple.

For ternary search in \(t\), start at \([0,M]\). While the interval has length greater than \(P^{-75}\), round its one-third and two-thirds positions to the nearest points of the \(P^{-95}\) grid to get \(t_1<t_2\) inside the interval. If the estimated value at \(t_1\) is at most that at \(t_2\), retain the interval through \(t_2\). For a point \(r>t_2\) in the discarded part, convexity gives
\[
 F(t_2)-F(r)\le
 (t_2-r)\frac{F(t_2)-F(t_1)}{t_2-t_1}
 \le 2\mu\frac{r-t_2}{t_2-t_1}.
\]
If the estimated value at \(t_1\) is larger, retain the interval starting at \(t_1\). For \(r<t_1\) in the discarded part, convexity gives
\[
 F(t_1)-F(r)\le
 (t_1-r)\frac{F(t_2)-F(t_1)}{t_2-t_1}
 \le 2\mu\frac{t_1-r}{t_2-t_1}.
\]
For interval length \(a>P^{-75}\), put \(h_t=P^{-95}\). The ratios are at most \((a/3+h_t)/(a/3-2h_t)<2\), and the new interval has length at most \(2a/3+h_t<3a/4\). Since \((3/4)^4<1/2\) and \(M=P^{15}\), this requires at most \(400\log_2 P+1\) comparisons. Each increases the minimum over the retained interval by at most \(4\mu\), and choosing its final left endpoint costs at most \(P^{-75}\) by the Lipschitz bound. On the good event the original minimum is at most \(P^{-96}\) by the retained-flow witness. Thus the final endpoint has value at most \(P^{-96}+4(400\log_2 P+1)P^{-80}+P^{-75}<P^{-56}\) using \(\log_2 P\le P\). Use a final tree call at tolerance \(P^{-56}\) in every component, with its chosen \(t\) in the unicyclic case. Reject the trial if any final tree call fails. The final maximum residual on a successful call is at most \(P^{-56}+(P+1)h_{\rm tree}<P^{-55}\). These search limits bound work on every trial. Summing the traversal sizes across components gives \(O((m_0+n_0)\operatorname{polylog} P)\) including all search calls and recovery.

Set other coordinates to zero. On the good event the recovered exact dyadic \(f\in[0,M]^{m_0}\) satisfies
\[
 (c')^\top f\le (b')^\top y+n_0(2P)P^{-55}+m_0M\theta
              \le (b')^\top y+P^{-7}.
\]
This is the identity \((c')^\top f=(b')^\top y+(L^\top f-b')^\top y+\sum_{e\in S} f_e(c'_e-L_e y)\) with \(\|y\|_\infty\le2P\), \(n_0(2P)P^{-55}\le2P^{-53}\) and \(m_0M\theta\le P^{-8}\). Check dual feasibility, the graph test, \(0\le f\le M\), perturbed residual at most \(P^{-55}\) in maximum norm, and \((c')^\top f\le(b')^\top y+P^{-7}\) by exact scans. Retry with fresh perturbations on a failed test. The good event implies acceptance, so the conditional failure probability per trial is at most \(10P^{-11}\). Every acceptance certifies the claimed box and perturbed residual, and weak duality certifies \((b')^\top y\le\mathrm{OPT}'\). The fresh trials give almost-sure termination and constant expected multiplicity. The probability of requiring more than \(j\) trials in a call is at most \((10P^{-11})^j\) conditional on the caller's history.

With respect to the unperturbed augmented problem let \(e=L^\top f-b\) and \(E_{\max}=\sigma+P^{-55}\), so \(\|e\|_\infty\le E_{\max}\). By (4) and the acceptance tests,
\[
 c^\top f\le(c')^\top f\le\mathrm{OPT}_{\rm box}+P^{-7}
                +2\sigma\sum_a U_a+n_0(H_a+2\sigma)\sigma
           \le\mathrm{OPT}_{\rm box}+2P^{-7}.
\]
Write \(t_i\) for the signed artificial contribution at vertex \(i\). Because each artificial unit has cost \(H_a\) and the remaining costs are nonnegative, the sum of artificial flows, and hence \(\|t\|_1\), is at most \((C+1)/H_a\). At the private vertex \(k=k(a)\) the equation reads \(x_a+z_a=U_a+t_k-e_k\). Substitute \(z_a=U_a-x_a+t_k-e_k\) in the old vertex \(i\)'s equation. After removing the artificial flows and subtracting the demand increments, the unperturbed original residual there, before clipping, is
\[
  e_i-t_i-\sum_{a:j(a)=i}\beta_a(t_{k(a)}-e_{k(a)}).
\]
Clip each \(x_a\) to \([0,U_a]\). Since \(x_a,z_a\ge0\), the amount clipped is at most \(|t_k|+|e_k|\), and its effect at an incident old vertex is bounded by that amount. Since the private vertices are distinct, the total per-old-vertex bound after clipping is at most \(3(C+1)/H_a+(1+2m_0)E_{\max}\). The first term is at most \(3\epsilon_{\rm out}/64\), and the second is at most \(6P^{-11}\) since \(m_0\le P\). Clipping and dropping auxiliary coordinates can only decrease the nonnegative objective. Thus the objective error is at most \(2P^{-7}\) and the equation error is at most \(3\epsilon_{\rm out}/64+6P^{-11}\), both at most \(\epsilon_{\rm out}\) by \(P\ge\epsilon_{\rm out}^{-1}\). The clipped box is exact.

For the center computation on either sign use the range bound \(V_{\rm ctr}=4m_0H^{7p}+32n_0p\log_2 H+4\). On its box \(Q\le\mathcal E\le m_0H^{7p}\) and \(|B(v_i)|\le6p\log_2H\). The \(O(p\log H)\) breakpoints per utility and \(O(p\log(Hp))\) per row give size \(O((m_0+n_0)p\log(Pp))\) for each model. There are \(O((p\log P)^2)\) models by the outer and inner iteration bounds, since \(\log V_{\rm ctr}=O(p\log P)\), the target is \(1/64\), and the outer loop uses a factor-two tolerance within a constant factor of the target divided by its round count. The \(T\) affine box problems on that sign use range \(V_{\rm aff}=4+4(\|g\|_1+M_bm_0H)\max_i D_i\), since each hinge is between zero and \(H\max_i D_i\). Its logarithm and the logarithm of the inverse error used there are \(O(\log P)\). Each affine box uses only the inner loop, giving \(O(\log P)\) models on \(O(m_0+n_0)\) arcs.

In the center inputs each breakpoint construction uses integers of bit length \(O(p(\log H+(J+1)\log(p+1)))=O(p^2\log^2 P)\), including the exact powers and the slope divisions. The unary breakpoints and the box endpoints use \(O(p\log P)\) bits. The affine inputs have \(O(\log P)\) bits because their grids use fixed powers of \(P\). All mesh choices and the two integer ranges of a model therefore satisfy \(\log(2+U_f+C_f)=O((1+p)^2\log^2 P)\). With \(R=O(\log^2 P)\), there are \(O((1+p)^2\log^4 P)\) exact ordinary-flow calls per trial, two per model across both signs. This gives \(O((1+p)^6\log^8 P)\) calls to the polynomial-range Lemma 2.1, by the explicit integer scaling count.

The row and anchor arrays, signed scores and physical centers are initialized or updated by scans. Center model construction uses polynomial bit work per generated term, including dyadic powers and divisions. Root extraction uses \(O(\log P)\) integer comparisons per coordinate on the \(P^{-12000}\) grid. Box and residual network construction scans the term lists on every rebuild. The ordinary solver bound includes all its queries, updates and rebuilding. Together with the recovery scans and searches, a trial costs \((m_0+n_0)^{1+o(1)}\operatorname{polylog}P\). To check the product form, put \(S_0=m_0+n_0\) and bound ordinary size by \(S_0D\), where \(D=\operatorname{polylog}P\) since \(p=O(\log P)\). For ordinary sizes at most \((S_0+2)^2\), the subpolynomial factor is bounded by \((S_0+2)^{o(1)}\). For larger sizes \(S_0+2=O(D)\), so each ordinary size is \(O(D^2)\) and its subpolynomial factor is bounded by a fixed polynomial in \(D\). This accounts for the same time bound at every graph density.

\hypertarget{arbitrary-signed-numerical-inputs-and-exact-original-cost-dominance}{%
\section*{Arbitrary signed numerical inputs and exact original cost dominance}\label{arbitrary-signed-numerical-inputs-and-exact-original-cost-dominance}}

Use a precision-access model with exact sign and zero metadata, exact comparisons of input cost magnitudes, and certified dyadic enclosures of normalized significands to \(k\) binary places. For a nonzero input \(x\), write \(|x|=2^{E(x)}s(x)\), with \(E(x)=\lfloor\log_2|x|\rfloor\) and \(1\le s(x)<2\). The enclosure for \(s(x)\) has length at most \(2^{-k}\) and endpoints of \(O(k+1)\) bits. A scale \(2^{E(x)}\) may be stored by reference to \(x\). An exponent query gives \(E(x)-E(y)\) clipped to \([-K,K]\), with a flag for overflow, also allowing \(y=1\). Charge polynomially many bit operations in \(k+\log(2+K)+\log N\) per precision query. Count sign and zero metadata access and exact cost comparisons as input operations on referenced data. For an explicitly supplied arbitrary-length encoding, charge the required reading, metadata extraction and exact input comparisons on that encoding separately. Charge expansion of a scale exponent or the full numeric expansion of an affine output separately when requested. Significand enclosure requests below use \(O(\log((N+2)(B+2)))\) leading bits, where \(B\) is defined next.

Discard zero capacities by a scan and record their positions for zero output. With no surviving arcs, feasibility gives \(d=0\), and the zero flow works. Let \(D_0=2^{E(\delta)}\), so \(D_0\le\delta<2D_0\), and let \(U_e=u_e/D_0\). Set \(j=\lceil\log_2m\rceil\). Select an arc with largest \(E(u_e)\) in one scan using clipped exponent comparisons with cutoff one. Query that exponent relative to \(E(\delta)\) with cutoff initially \(j+2\), doubling the cutoff on positive overflow. On negative overflow use the negative cutoff. With \(h_u\) the resulting difference, set \(B=2^{\max(0,1+j+h_u)}\). Without negative overflow, \(2^{h_u}\le\max_e U_e<2^{h_u+1}\), so \(B\) is between \(\max(1,m\max_e U_e)\) and four times this value by \(m\le 2^j<2m\). On negative overflow \(B=1\) and \(m\max_e U_e<1\). The final cutoff is \(O(1+\log m+\log B)\). Thus \(B\le4\max(1,2m\max_e u_e/\delta)\) and \(\log(2+B)=O(\log(2+NW/\delta))\) for every \(\delta>0\). Let \(f^0_e=u_e\) at negative costs and zero otherwise. With \(\chi_e=-1\) at negative costs and \(+1\) otherwise, write
\[
 f=f^0+D_0\operatorname{diag}(\chi)y,\quad
 A=A_0\operatorname{diag}(\chi),\quad b=(d-A_0 f^0)/D_0,\quad
 a_e=|c_e|.
\]
The affine formula bijects \(0\le y\le U\) with \(0\le f\le u\), and \(A_0f-d=D_0(Ay-b)\). Also \(c_e\chi_e=a_e\), so \(c^\top f=c^\top f^0+D_0a^\top y\). Therefore the true normalized optimal penalty is \(p_*=\min\{a^\top y:Ay=b,\ 0\le y\le U\}\ge0\) and \(\mathrm{OPT}=c^\top f^0+D_0p_*\). Each entry of \(A\) and \(A_0\) has magnitude at most one, including the combined entry of a loop. Every normalized box point has \(\|y\|_1\le\sum_e U_e\le B\). Feasibility of the exact normalized problem gives \(\|b\|_\infty\le B\). Original feasibility also gives \(\|d/D_0\|_\infty\le B\), since \(d/D_0=A_0(x/D_0)\) for some \(0\le x\le u\).

For a normalized input \(x/2^{E(y)}\) of magnitude at most \(R\), with \(R\ge1\), a target dyadic mesh \(0<h<1\), and \(x\ne0\), query the exponent difference with \(K=\max(\lceil\log_2 R\rceil+2,\lceil\log_2(16/h)\rceil+2)\). Negative overflow gives an enclosure on the known sign side of zero of width at most \(2^{1-K}\le h/4\). Positive overflow is excluded by the magnitude bound. Otherwise query \(k=\lceil\log_2(16R/h)\rceil\) leading places of \(s(x)\), and multiply the enclosure by the returned power of two, with the input sign. Its width is at most \(R2^{-k}\le h/16\). A zero input has the singleton enclosure. Round using these enclosures as follows, with \(\eta_0=1/1024\) and \(j=\lceil\log_2m\rceil\).

\begin{itemize}
\tightlist
\item
  For \(U_e=u_e/D_0\), use magnitude bound \(B\) and mesh \(h_U=\eta_0/(64\cdot 2^j)\). Round the lower enclosure endpoint down and clip at zero. Then \(U'_e\in[\max(0,U_e-\eta_0/(8m)),U_e]\), since the combined error is at most \((1+1/4)h_U\). Omit variables with \(U'_e=0\), restoring their zero normalized coordinates on output.
\item
  For \(\gamma_e\), use reference \(1\), magnitude bound one and mesh \(h_\gamma=\eta_0/(16B)\). Round the enclosure midpoint to the mesh and clip to \([0,1]\), at absolute error at most \(\eta_0/(8B)\). Form \(A'\) using the original column signs, combining loops and omitting zero terms.
\item
  For \(d_v/D_0\), use bound \(B\) and mesh \(h_d=\eta_0/32\). Rounding the midpoint to the mesh incurs absolute error at most \(\eta_0/4\). To compute \(b'\), subtract the contributions of the negative-cost origins using \(U'\) and the rounded gains, in normalized units.
\end{itemize}

View \(A'\) on all original positive-capacity columns for comparisons. A variable omitted or fixed in an LP has normalized coordinate zero in the corresponding output. Replacing the widths in the origins incurs at most \(\eta_0/8\) per row because the sum of absolute width changes is at most that amount and every coefficient has magnitude at most one. Replacing the gains costs at most \((\eta_0/(8B))\sum_e U_e\le\eta_0/8\) per row, including for loops. Thus \(\|b-b'\|_\infty\le\eta_0\) and \(\|(A-A')y\|_\infty\le\eta_0\) on the original box. Dyadic magnitudes and denominators, and the reciprocals of nonzero input coefficients of the rounded problem, are polynomial in \(N+B\). Truncating an exact normalized optimizer to \(U'\) changes \(A'y\) by at most the sum of width errors. It therefore gives a point of residual at most \(3\eta_0\) in \(A'y=b'\) and true penalty at most \(p_*\), using nonnegative \(a\). Every point of the smaller box has true normalized residual at most its rounded residual plus \(2\eta_0\).

To identify a positive cost unit, test allowed prefixes of the columns. Give all allowed physical variables zero objective coefficient in the test LP and fix the other physical variables to zero. Append a zero-cost signed unit slack in \([-1/16,1/16]\) per vertex and two opposite unit slack variables of cost one in \([0,B+3]\). Enforce the rounded demands \(b'\). This LP always has a box witness with zero physical flow by assigning the positive and negative parts of \(b'\) to the cost-one slacks, since \(\|b'\|_\infty\le B+\eta_0\). Solve to additive accuracy \(\eta_0\) by the bounded solver. Accept if its returned physical \(y\) has \(\|A'y-b'\|_\infty\le1/16+2\eta_0\). Whenever an allowed point has residual at most \(1/16\), the free slacks alone can supply the residual, so the test optimum is zero. The total cost-one slack returned is then at most \(\eta_0\) and the solver's equation error is at most \(\eta_0\). Along with the box bound on the free slack this proves acceptance in that case. Each acceptance supplies the physical witness verified by the residual check.

First permit only zero-cost arcs. On acceptance output its affine flow directly. It lies in the original box, has zero penalty, and has true normalized residual at most \(1/16+4\eta_0<1\), so its original imbalance error is at most \(D_0\le\delta\). Otherwise there is a positive cost magnitude, because allowing all arcs permits the truncated optimizer of residual at most \(3\eta_0<1/16\). Sort the distinct positive input cost magnitudes and allow increasing prefixes, retaining the zero-cost arcs in every prefix. Test the all-permitted prefix to initialize the accepted endpoint, using the zero-cost-only prefix as the rejected endpoint. Use bisection by replacing the upper endpoint on acceptance and the lower endpoint on rejection, always keeping the accepted witness. Each test reduces the length of the index interval, so the indices become adjacent even if decisions are ambiguous near the tolerance. Write \(\lambda>0\) for the magnitude allowed at the upper index and excluded at the lower. If \(p_*/\lambda\le1/16-3\eta_0\), deleting all coordinates whose cost magnitude is at least \(\lambda\) from the truncated optimum loses at most \(p_*/\lambda\) total flow. Each matrix coefficient has magnitude at most one, so the remaining point has residual at most \(1/16\). That would force acceptance at the adjacent rejected index. Thus \(p_*/\lambda>1/16-3\eta_0>1/32\). Every positive coordinate of the accepted witness has magnitude at most \(\lambda\) in its cost coefficient, so its penalty is at most \(B\lambda\).

Set \(t=1/(16B),\ \xi=t/1024,\ \zeta=\xi/B\). Approximate \(\min(a_e/\lambda,128B)\) by a nonnegative dyadic \(q_e\) to absolute error \(\zeta\). Handle a zero numerator directly. For positive magnitudes the quotient of normalized significands lies between \(1/2\) and 2. Exponent differences at least \(1+\log_2(128B)\) give the cap, and differences at most \(\lfloor\log_2\zeta\rfloor-2\) permit zero. Use a clipped exponent query covering these two thresholds. For intermediate differences query the two significands to \(\lceil\log_2(2^{14}B/\zeta)\rceil\) places, with enclosures intersected with \([1,2]\). Dividing the lower numerator by the upper denominator and conversely encloses the quotient to width at most six times the significand mesh, since the denominators are at least one and the numerators at most two. Scaling by the exponent difference multiplies that width by at most \(256B\). Clip the enclosure to the capped range by applying the cap to its endpoints and round its midpoint on the \(\zeta/2\) grid. This gives the claimed absolute error with \(O(\log(2+B))\) leading places. Remove variables with \(q_e>64B\). Each has \(a_e/\lambda>63B\), and each retained coefficient has \(|q_e-a_e/\lambda|\le\zeta\) since a capped ratio would have been removed. The accepted witness uses only magnitudes at most \(\lambda\), so all its positive coordinates are retained.

A comparator in this retained box has rounded residual at most \(1/8\) and true penalty at most \(p_*\). For \(p_*\ge B\lambda\) use the accepted witness. For \(p_*< B\lambda\) delete the removed entries of the truncated optimum, whose total mass is less than \(1/63\), giving residual at most \(3\eta_0+1/63\). Solve the retained rounded LP with one free signed slack in \([-1/8,1/8]\) per vertex, objective \(q^\top y\), and accuracy \(\xi\). It is exactly feasible by the comparator. Using \(\|y\|_1\le B\) for both the returned point and the comparator gives
\[
 \frac{a^\top y}{\lambda}\le\frac{p_*}{\lambda}+2B\zeta+\xi
                     =\frac{p_*}{\lambda}+3\xi,\qquad
 \|A'y-b'\|_\infty\le1/8+\xi .
\]
Return the affine flow with normalized physical coordinates \(\tilde y=(1-t)y\). Since \(3\xi<t/32<t p_*/\lambda\), the penalty bound after division by \(\lambda\) is \((1-t)(p_*/\lambda+3\xi)\le p_*/\lambda-tp_*/\lambda+3\xi\le p_*/\lambda\). Thus the original cost is at most \(\mathrm{OPT}=c^\top f^0+D_0p_*\). Its true normalized demand error is at most \(1/8+\xi+2\eta_0+tB<1\), since \(\|Ay\|_\infty\le B\). In original units the error is at most \(D_0\le\delta\). Since \(0\le\tilde y\le U'\le U\), the affine formula gives exact original box bounds. Save the normalized offsets as dyadics with the scale reference for \(D_0\). At negative cost the representation is \(u_e-D_0\tilde y_e\) with a reference to the exact \(u_e\). At other costs it is \(D_0\tilde y_e\). The zero-cost-prefix case uses the same representation with its returned normalized witness in place of \(\tilde y\).

There are \(J=O(\log(2+m))\) bounded solver calls. Sorting uses \(O(N\log(2+N))\) exact input cost comparisons. Arrays store the input indices, signs, rounded columns and prefix endpoints. The accepted witness is stored when a new accepted endpoint is found. Computing the scale uses \(O(N+\log(2+\log B))\) exponent queries. Every cutoff at this stage is \(O(1+\log N+\log B)\), and the same bound covers all cutoff arguments in the later normalization and cost-ratio queries. The representative input cost for \(\lambda\) is stored by its index. Subsequent comparisons involving computed values use only the bounded dyadics and rationals specified above. All other preprocessing, exact witness checks, prefix construction and output arithmetic take \(O(N\operatorname{polylog}(N+B+2))\) time across the calls. The augmented LPs have \(O(N)\) rows and columns, with all magnitude bounds, denominators, nonzero reciprocals and inverse target errors at most \((N+B+2)^{O(1)}\). Thus their choices of \(P\ge N\) have \(\log P=O(\log(N+B+2))\). Output coordinates use the tree grid, the cycle parameter grid or a dyadic width endpoint. Multiplication by \(1-t\) still gives \(O(\log(N+B+2))\) bits per normalized offset, plus the indices of the exact input data referenced by the output.

By the conditional failure bound and a union bound, the probability that any reached call needs a second trial is at most \(10J N^{-11}\). For any larger fixed high-probability exponent, allow a correspondingly larger fixed constant number of trials per call in the time bound, using the conditional geometric tail. Hence with high probability the end-to-end time is \(N^{1+o(1)}\operatorname{polylog}(N+B+2)\), with the same bound in expectation. Factors polynomial in \(\log N\) are absorbed into \(N^{o(1)}\), and \(\log(2+B)=O(\log(2+NW/\delta))\). All guarantees hold on every accepted output, and the time analysis includes the sparse case.

\hypertarget{relation-to-the-specified-references}{%
\section*{Relation to the specified references}\label{relation-to-the-specified-references}}

The proof lifts the capacity conversion of reference 2, Appendix B, to opposite-sign two-entry columns and carries the slack scaling of its dual \(\ell_1\) IPM into the displacement model. The piecewise-linear center regularizer, affine allowance penalties and seed mixture give the coordinate bounds and scale recourse proved above. The recourse bookkeeping adapts reference 3, Sections 5.5--5.6, and the ordinary subproblems use its Lemma 2.1, based on the \(\ell_1\) IPM and min-ratio-cycle data structure for ordinary flow. \href{https://arxiv.org/pdf/2407.10830v1}{Reference 2} \href{https://github.com/openai/math/blob/main/preprints/Almost-Linear-Time-Maximum-Cardinality-Matching-in-Sparse-General-Graphs-September-24-2026/main.pdf}{Reference 3, September 24, 2026}

For context, Jiang, Kapralov, Li and Sidford, Theorem 1.4 and Section 8.4, give an additive objective guarantee with an \(n^{1.5}\) term in the running time. Daitch and Spielman, Section 3.3 and Lemma 3.9, use cost-shrinking for min-cost maximum flow with positive integral costs. Here the accepted and rejected prefix tests supply a lower bound in the chosen cost unit, which makes the shrinking argument apply to a signed real-cost instance with prescribed demands. \href{https://arxiv.org/pdf/2510.17740}{Reference 1, Theorem 1.4 and Section 8.4} \href{https://arxiv.org/pdf/0803.0988}{Reference 4, Section 3.3 and Lemma 3.9}

\end{document}
